%%%%%
%
\section{Rotation horizontale appliqu�e � un vecteur vent pour passer dans le rep�re g�ographique}

\itemstyle{Op\'erateur~:} {\sf Rotat\_uv}
\index{Rotat\_uv}

%%%%%OBJET%%%%
\itemstyle{Objet~:} 
Applique une rotation horizontale � un vecteur (u,v) exprim� dans un rep�re orthogonal [$\overrightarrow{i},\overrightarrow{j},\overrightarrow{k}$] pour passer dans un rep�re g�ographique avec les conventions dites "m�t�o".

Cet op�rateur permet de traiter 2 cas : le rep�re direct peut �tre soit  [$\overrightarrow{i},\overrightarrow{j},\overrightarrow{k}$], soit
[$-\overrightarrow{j},\overrightarrow{i},\overrightarrow{k}$].

%%%%%ARGUMENTS%%%%
\bigskip
%%%%%ENTREES%%%%
\begin{tabarg}
\multicolumn{3}{l}{\titretab {Entr\'ees}}\\
\hline
$U$ & Vecteur, {\tt float} & Premi�re composante horizontale (unit� quelconque)\\
$V$ & Vecteur, {\tt float} & Seconde composante horizontale (m�me unit� que $U$)\\
$offset$ & Coeff., {\tt float} & Angle (sens trigonom�trique) entre le Nord g�ographique et U (en degr�s)\\
$flag$ & Coeff., {\tt float} & Flag � 1 (i,j,k) direct, ou -1 (-j,i,k) direct (sans unit�)\\
\end{tabarg}

\medskip
%%%%%SORTIES%%%%
\begin{tabarg}
\multicolumn{3}{l}{\titretab {Sorties}}\\
\hline
$U_{G}$ & Vecteur, {\tt float} & Composante Est, compt�e depuis l'Ouest (m�me unit� que $U$ et $V$)\\
$V_{G}$ & Vecteur, {\tt float} & Composante Nord, compt�e depuis le sud (m�me unit� que $U$ et $V$)\\
\end{tabarg}



%%%%%FORMULE%%%%
\itemstyle{Formule}

%%Rotat_uv

On ajoute $\pi$/2 � l'offset qui est donn� par rapport au Nord et non � l'Est :
$new\_offset  =  \frac{\pi}{2} + offset . \frac{\pi}{180}$
On pose :
\begin{align*}
coso &= \cos (new\_offset)\\
sino &= \sin (new\_offset)
\end{align*}
On en d�duit les nouvelles composantes not�es $U_{G}$ et $V_{G}$ :
\begin{align*}
U_{G} &= flag . coso . U  - sino . V\\
V_{G} &= flag . sino . U  + coso . V
\end{align*}




%%%%%COMMENTAIRE%%%%
\itemstyle{Commentaire~:}\\
\underline{Application � des donn�es de vents d'un an�mom�tre sonique comme le GILL :}\\
L'offset est la somme de l'angle d'orientation du GILL par rapport au Nord g�ographique (ou parfois l'avant du bateau) et de l'angle (30 degr�s sur un GILL HS) entre l'axe support de U et l'axe de l'armature du GILL ("Nord du GILL").\\
\textbf{Attention}, tout est compt� dans le sens trigonom�trique.\\
%%%%%AUTEUR%%%%
\itemstyle{Auteur~:}
CNRM/GMEI/TRAMM

%%%%%REFERENCE%%%%
\itemstyle{R\'ef\'erences~:}

